
\def\PUDcodigo{ECA209}

\if\COMPILAR0
   \def\PUDcodacad{04507.14}
\else
   \def\PUDcodacad{04506.15}
\fi

\def\PUDdisciplina{Linguagem de Programação}

\def\PUDsemestre{S3}

\def\PUDhoras{80}

%NOVOS ---------------------------------------------------
\def\PUDhorasTeoria{40}
\def\PUDhorasPratica{40}

\def\PUDinterECA{ 
\hyperref[PUD:ECA208]{Lógica de Programação (04507.10)}

\hyperref[PUD:ECA211]{Eletrônica I (04507.15)}
}

\def\PUDatuaisECA{
- Aplicações com tecnologias atuais
}

\def\PUDrecursos{Projetor multimídia; Quadro branco e pincel; Aulas em laboratório de informática.}

%----------------------------------------------------------

\def\PUDcreditos{4}

\def\PUDprerequisito{ \hyperref[PUD:ECA208]{ECA208} }

\if\COMPILAR0
   \def\PUDpreacad{ \hyperref[PUD:ECA208]{Lógica de Programação (04507.10)} }
\else
   \def\PUDpreacad{ \hyperref[PUD:ECA208]{Lógica de Programação (04506.11)} }
\fi

\def\PUDobjetivos{Compreender o conceito de abstração de dados, sua importância para os princípios de modularidade, encapsulamento e independência de implementação. Ler e escrever programas de computador orien- tados a objetos. Compreender as estruturas de dados clássicas, suas características funcionais, formas de representação, operações associadas e complexidade das operações. Além de verificar o contexto de aplicações regionais e interdiciplinariedade tanto com as disciplinas de pré-requisito quanto com as disicplinas do semestre corrente, como Eletrônica I e outras.
%Compreender o conceito de abstração de dados, sua importância para os princípios de modularidade, encapsulamento e independência de implementação.  Ler e escrever programas de computador orientados a objetos.  Compreender as estruturas de dados clássicas, suas características funcionais, formas de representação, operações associadas e complexidade das operações. Além de verificar o contexto de aplicações regionais e interdiciplinariedade tanto com as disciplinas de pré-requisito quanto com as disicplinas do semestre corrente, como Eletrônica I e outras.
} %A ÚLTIMA FRASE É NOVA.

\def\PUDementa{Introdução à programação orientada à objetos através do uso da Linguagem C/C++. Ponteiros. Pilhas. Filas. Listas Ordenadas. Técnicas Avançadas de Encadeamento. Recursividade. Lista Lineares. Listas Generalizadas. Ordenação de Dados.
%Programas Procedimentais x Programas Orientados a Objetos.  Lista Lineares.  Introdução a programação orientada a objetos.  Alocação de Memória Estática, Dinâmica, Seqüencial e Encadeada.  Pilhas.  Filas.  Listas Ordenadas.  Técnicas Avançadas de Encadeamento.  Recursividade.  Listas Generalizadas.  Árvores.  Ordenação de Dados.  Linguagem adotada C/C++.
}

\def\PUDconteudo{ & Conceitos de programação orientada a objetos. Programas Procedimentais x Programas Orientados a Objetos. Objetos e Classes. Herança e Polimorfismo. Encapsulação. Agregação e Composição. Interfaces.
\\
& Listas lineares. Definição e operações aplicáveis. Implementação utilizando vetor. 
\\
& Tipos de implementação.  Alocação de Memória Estática e Dinâmica.  Alocação de
Memória Sequencial e Encadeada. Utilização de ponteiros.
\\
& Pilhas. Definição e operações aplicáveis. Implementação. 
\\
& Filas. Definição e operações aplicáveis. Implementação .Aplicações clássicas
\\
& Listas ordenadas. Definição e operações aplicáveis. Implementação. Aplicações clássicas: Mapeamentos, Polinômios e Filas de Prioridade.
\\
& Recursividade. Conceito de recursividade. Sequências definidas recursivamente. Operações definidas recursivamente.
\\
& Listas generalizadas. Definição e operações aplicáveis. Implementação utilizando vetores.
%& Conceitos de programação orientada a objetos. Programas Procedimentais x Programas Orientados a Objetos. Objetos e Classes. Herança e Polimorfismo. Encapsulação. Agregação e Composição. Interfaces. 
%\\ 
% &  Listas lineares. Definição e operações aplicáveis. Implementação utilizando vetor. 
%\\ 
% &  Tipos de implementação. Alocação de Memória Estática e Dinâmica. Alocação de Memória Seqüencial e Encadeada. 
%\\ 
% &  Pilhas. Definição e operações aplicáveis. Implementação.  Aplicação clássica: Avaliação de expressões. 
%\\ 
% &  Filas. Definição e operações aplicáveis. ImplementaçãoAplicação clássica: Colorindo regiões gráficas. 
%\\ 
% & Listas ordenadas. Definição e operações aplicáveis. Implementação. Aplicações clássicas: Mapeamentos, Polinômios e Filas de Prioridade. 
%\\ 
% & Técnicas avançadas de encadeamento. Nodos cabeça e sentinela. Encadeamento circular. Encadeamento duplo. Encadeamento duplo compactado. 
%\\ 
% & Recursividade. Conceito de recursividade. Seqüências definidas recursivamente. Operações definidas recursivamente. 
%\\ 
% & Listas generalizadas. Definição e operações aplicáveis. Implementação
%\\ 
% & Árvires. Conceitos sobre árvore. Árvore binária. Árvore de busca binária. Implementação de árvore de busca binária. Aplicação clássica: Compactação de dados. 
%\\ 
% & Ordenação de dados. Ordenação por inserção. Ordenação por troca. Ordenação por seleção. 
%\\ 
% & Comparação entre os métodos. Eficiência de algoritmos: A notação Big-O. 
%\\ 
% & PEsquisa de dados. Pesquisa seqüencial. Pesquisa binária. 
}

\def\PUDmetodo{Aulas expositórias que podem ser teóricas e/ou práticas, onde as práticas no laboratório serão marcadas no decorrer da disciplina.}

\def\PUDavalia{A avaliação será feita com aplicação de provas teóricas e/ou práticas, além da possibilidade de inclusão de trabalhos e seminários no decorrer da disciplina.}

\def\PUDbibbasica{ &  \href{www.biblioteca.ifce.edu.br/index.asp?codigo_sophia=24506}{Deitel}, H. M.; Deitel, P. J.  C++: como programar.  6º ed.  São Paulo, SP: Pearson Prentice Hall, 2011.  818p.  ISBN: 9788576059349.
\\ 
 &  \href{www.biblioteca.ifce.edu.br/index.asp?codigo_sophia=40651}{Forbellone}, André Luiz Villar.  Lógica de programação: a construção de algoritmos e estruturas de dados.  3º ed.  São Paulo, SP: Pearson, 2013.  218p.  ISBN: 9788576050247.
%\\ 
 %&  IFAO - INFORMATIONSSYSTEME GMBH.  Comando numérico CNC: técnica operacional: curso básico.  São Paulo (SP): EPU, 1984.  176 p.  ISBN 9788512180700
\\ 
 & \href{www.biblioteca.ifce.edu.br/index.asp?codigo_sophia=24502}{Souza}, M. A. F.; Gomes, M. M.; Soares, M. V. Concilio, R.  Algoritmos e Lógica de Programação.  São Paulo, SP: Cengage Learning, 2008.  214p.  ISBN: 8522104646. }

\def\PUDbibcomplementar{ & \href{www.biblioteca.ifce.edu.br/index.asp?codigo_sophia=62925}{Ziviani}, Nivio.  Projeto de algoritmos: com implementações em Java e C++.  São Paulo, SP: Cengage Learning, 2015.  621p.  ISBN: 9788522105250.
\\ 
 & \href{www.biblioteca.ifce.edu.br/index.asp?codigo_sophia=24429}{Beneduzzi}, Humberto Martins.  Lógica e linguagem de programação: introdução ao desenvolvimento de software.  144p.  ISBN: 9788563687111.
%\\ 
 %& KERNIGHAN, B.  W. , RITCHIE, D.  M. , C: A linguagem de programação.  Rio de Janeiro: Elsevier, 1986.  208p.  ISBN 9788570015860. 
\\ 
 & \href{www.biblioteca.ifce.edu.br/index.asp?codigo_sophia=65446}{Cormen}, T. H.  Algoritmos: teoria e prática.  3º ed.  Rio de Janeiro, RJ: Elsevier, 2012.  926p.  ISBN: 9788535236996.
\\
 & \href{www.biblioteca.ifce.edu.br/index.asp?codigo_sophia=24241}{Chapman}, Stephen J.  Programação em MATLAB para engenheiros.  2º ed.  São Paulo, SP: Cengage Learning, 2010.  410p.  ISBN: 9788522107896.
\\
 & \href{www.biblioteca.ifce.edu.br/index.asp?codigo_sophia=23797}{Hanselman}, Duane.  MATLAB 6: curso completo.  São Paulo, SP: Prentice Hall, 2003.  676p.  ISBN: 9788587918567. }

\def\PUDautor{Geraldo Luis Bezerra Ramalho}

\def\PUDdata{2013-04-17}

\def\PUDrevisao{3}

\def\PUDrevisor{Antonio Barbosa de Souza Júnior}

\def\PUDdatarevisao{2019-05-15}

\PUDcorpo
